% Source: https://artofproblemsolving.com/community/c3h1549767_more_questions_relating_to_gcd_lcm
% Stable identity: aops-topic-1549767-q04
Let $g=\gcd(a,b)$ and write $a=gx$, $b=gy$ with $\gcd(x,y)=1$. Then $\operatorname{lcm}(a,b)=gxy=1764=2^2\cdot3^2\cdot7^2$ and $g^2(x^2+y^2)=85113=3^2\cdot7^2\cdot193$.

So $g^2\mid3^2\cdot7^2$, i.e.\ $g\in\{1,3,7,21\}$, and $193$ is prime. Try each:
\begin{itemize}
\item $g=21$: $x^2+y^2=193$ and $xy=84$. Then $(x+y)^2=193+168=361$ and $(x-y)^2=193-168=25$, so $x=12$, $y=7$.
\item $g=7$: $x^2+y^2=1737$, $xy=252$: $(x+y)^2=2241$, not a square.
\item $g=3$: $x^2+y^2=9457$, $xy=588$: $(x+y)^2=10633$, not a square.
\item $g=1$: $x^2+y^2=85113$, $xy=1764$: $(x+y)^2=88641$, not a square ($297^2=88209$, $298^2=88804$).
\end{itemize}
Hence $a=21\cdot12=252$, $b=21\cdot7=147$. Check: $252^2+147^2=63504+21609=85113$, and $\operatorname{lcm}(252,147)=\dfrac{252\cdot147}{21}=1764$. The answer is $\boxed{(252,147)}$.
